\documentclass[11pt]{article}

\usepackage[margin=1in]{geometry}
\usepackage[T1]{fontenc}
\usepackage[utf8]{inputenc}
\usepackage{lmodern}
\usepackage{hyperref}
\usepackage{xcolor}
\usepackage{listings}
\usepackage{enumitem}
\usepackage{longtable}

\hypersetup{
  colorlinks=true,
  linkcolor=blue,
  urlcolor=blue,
  pdftitle={TP CAP Replication Technical Document}
}

\lstdefinestyle{javaStyle}{
  language=Java,
  basicstyle=\ttfamily\small,
  keywordstyle=\color{blue}\bfseries,
  commentstyle=\color{gray},
  stringstyle=\color{teal!70!black},
  numbers=left,
  numberstyle=\tiny\color{gray},
  stepnumber=1,
  numbersep=8pt,
  frame=single,
  breaklines=true,
  tabsize=2,
  showstringspaces=false,
  columns=fullflexible
}

\lstset{style=javaStyle}

\title{TP CAP Replication with RabbitMQ}
\author{Technical Documentation}
\date{\today}

\begin{document}

\maketitle

\begin{abstract}
This project implements a small distributed replication system using Spring Boot and RabbitMQ. Three independent replicas store data in local files, a writer broadcasts updates to every replica, and two reader modes demonstrate the trade-off between availability and consistency. This document explains the architecture, the message flow, the storage model, the client behavior, and the failure scenarios directly reflected in the source code.
\end{abstract}

\tableofcontents

\section{System Overview}

The application is organized as a single Spring Boot jar that can run in multiple modes through Spring profiles. The same codebase starts as one of four roles:

\begin{itemize}[leftmargin=1.4em]
  \item \textbf{replica1, replica2, replica3}: listen for write messages and serve read requests.
  \item \textbf{writer}: send new lines to the replication fanout exchange.
  \item \textbf{reader}: request the last line from the first replica that answers.
  \item \textbf{readerv2}: request all lines from every replica and reconcile them using majority voting.
\end{itemize}

The design is intentionally simple: all persistent state is stored as plain text files under \texttt{replica\_data/}. This makes the replication behavior easy to inspect and also makes divergence visible when one replica misses a write.

\section{Build and Runtime Model}

The main entry point is the standard Spring Boot launcher:

\begin{lstlisting}
package com.tp.replication;

import org.springframework.boot.SpringApplication;
import org.springframework.boot.autoconfigure.SpringBootApplication;

@SpringBootApplication
public class TpCapReplicationApplication {
    public static void main(String[] args) {
        SpringApplication.run(TpCapReplicationApplication.class, args);
    }
}
\end{lstlisting}

There is no custom bootstrap logic here. All behavior is attached through profile-scoped beans, RabbitMQ listeners, and \texttt{CommandLineRunner} implementations. That means the same jar can play very different roles depending on the active profile passed at startup.

\section{Configuration and Profiles}

The runtime configuration lives in \texttt{application.yml}. The root profile defines the shared defaults, and each replica profile injects its own identity, data directory, and queue names.

\begin{lstlisting}
spring:
  application:
    name: tp-cap-replication
  rabbitmq:
    host: localhost
    port: 5672
    username: guest
    password: guest
  profiles:
    active: replica1

replica:
  count: 3

client:
  read-timeout-ms: 5000
  read-all-timeout-ms: 5000
\end{lstlisting}

Each replica profile adds the values needed by \texttt{ReplicaService}:

\begin{lstlisting}
replica:
  id: replica1
  data-dir: replica_data/replica1
  write-queue: replica_write_replica1
  read-queue: replica_read_replica1
\end{lstlisting}

This profile-based design is important because it avoids duplicating code across three separate applications. Instead, one service class reads its queue and file settings from the active profile and behaves like a distinct replica.

\section{RabbitMQ Topology}

RabbitMQ is used as the transport layer between clients and replicas. The configuration defines two fanout exchanges and separate replica queues for writes and reads.

\begin{lstlisting}
@Configuration
public class RabbitMQConfig {

    public static final String WRITE_EXCHANGE = "write_exchange";
    public static final String READ_REQUEST_EXCHANGE = "read_request_exchange";

    @Bean
    public FanoutExchange writeExchange() {
        return new FanoutExchange(WRITE_EXCHANGE, true, false);
    }

    @Bean
    public FanoutExchange readRequestExchange() {
        return new FanoutExchange(READ_REQUEST_EXCHANGE, true, false);
    }
\end{lstlisting}

The important architectural choice is the use of a fanout exchange. A fanout exchange delivers every message to every bound queue, which is exactly what a replication demo needs for write propagation and broadcast reads.

The configuration then creates one write queue and one read queue per replica and binds them to the exchanges:

\begin{lstlisting}
@Bean
public Queue replicaWriteQueue1() {
    return new Queue(REPLICA1_WRITE_QUEUE, false, false, true);
}

@Bean
public Binding replicaQueueBinding1(Queue replicaWriteQueue1, FanoutExchange writeExchange) {
    return BindingBuilder.bind(replicaWriteQueue1).to(writeExchange);
}
\end{lstlisting}

The queue flags mean the queues are non-durable, non-exclusive, and auto-delete. In practice this keeps the demo lightweight and disposable: when a replica process disappears, its queues do not need to persist as durable broker state.

\section{Local Persistence Model}

The replicas do not use a database. Instead, each replica appends to a single text file. That decision keeps the system easy to understand and creates a visible consistency story: if a replica misses a message, its file diverges from the others.

\begin{lstlisting}
public static void appendLine(String filePath, String line) {
    try {
        Path path = Paths.get(filePath);
        Files.createDirectories(path.getParent());

        String content = line + System.lineSeparator();
        Files.write(path, content.getBytes(StandardCharsets.UTF_8),
                StandardOpenOption.CREATE, StandardOpenOption.APPEND);
    } catch (IOException e) {
        logger.error("Failed to append to {}: {}", filePath, e);
    }
}
\end{lstlisting}

This method is the basic write primitive used by every replica. It creates the parent directory if needed, appends a newline-terminated record, and logs failures instead of crashing the process.

The read helpers are equally direct:

\begin{lstlisting}
public static String readLastLine(String filePath) {
    try {
        Path path = Paths.get(filePath);
        if (!Files.exists(path)) {
            return "";
        }

        List<String> lines = Files.readAllLines(path, StandardCharsets.UTF_8);
        if (lines.isEmpty()) {
            return "";
        }
        return lines.get(lines.size() - 1);
    } catch (IOException e) {
        logger.error("Failed to read last line from {}: {}", filePath, e);
        return "";
    }
}
\end{lstlisting}

The file helper returns an empty string when the file is missing or empty. That behavior keeps the reader clients simple because they can treat the absence of data as a valid, handled case.

\section{Replica Service}

\texttt{ReplicaService} is the core server-side component. It is only active under the replica profiles, and it handles both writes and read requests.

\subsection{Initialization}

On startup, the service computes the local data file path and ensures the replica directory exists.

\begin{lstlisting}
@PostConstruct
public void init() {
    if (replicaId != null && !replicaId.isEmpty()) {
        dataFilePath = dataDir + "/data.txt";
        FileUtils.ensureDirectory(dataDir);
        logger.info("Replica {} initialized with data directory: {}", replicaId, dataDir);
    }
}
\end{lstlisting}

This is a small but important step. The data path is not hardcoded; it comes from the active profile so each process writes to a different file tree.

\subsection{Write Path}

The write listener receives plain text lines from the \texttt{write\_exchange} fanout exchange.

\begin{lstlisting}
@RabbitListener(queues = "${replica.write-queue}")
public void handleWrite(String message) {
    logger.info("[{}] Received write message: {}", replicaId, message);
    FileUtils.appendLine(dataFilePath, message);
    logger.info("[{}] Data persisted to {}", replicaId, dataFilePath);
}
\end{lstlisting}

The method is intentionally minimal: every replica performs the same append locally. There is no consensus protocol and no coordination between replicas. The broker delivers the message, and each replica persists it independently.

\subsection{Read Request Parsing}

Read requests arrive on the replica-specific read queue and are parsed using a simple colon-delimited protocol.

\begin{lstlisting}
@RabbitListener(queues = "${replica.read-queue}")
public void handleReadRequest(String message) {
    String[] parts = message.split(":", 3);
    if (parts.length < 3) {
        logger.warn("[{}] Invalid read request format: {}", replicaId, message);
        return;
    }

    String requestType = parts[0];
    String correlationId = parts[1];
    String replyQueueName = parts[2];

    if ("READ_LAST".equals(requestType)) {
        handleReadLast(correlationId, replyQueueName);
    } else if ("READ_ALL".equals(requestType)) {
        handleReadAll(correlationId, replyQueueName);
    }
}
\end{lstlisting}

This protocol makes the request structure explicit:

\begin{itemize}[leftmargin=1.4em]
  \item \texttt{READ\_LAST:<correlationId>:<replyQueueName>}
  \item \texttt{READ\_ALL:<correlationId>:<replyQueueName>}
\end{itemize}

The correlation identifier allows the client to ignore stray replies from older requests, and the reply queue name tells replicas where to send the answer.

\subsection{READ\_LAST Scenario}

The availability-oriented read path returns only the last line from one replica.

\begin{lstlisting}
private void handleReadLast(String correlationId, String replyQueueName) {
    String lastLine = FileUtils.readLastLine(dataFilePath);
    String response = correlationId + "|" + lastLine;

    rabbitTemplate.convertAndSend(replyQueueName, response);
    logger.info("[{}] Sent READ_LAST response to {}: {}", replicaId, replyQueueName, response);
}
\end{lstlisting}

The response format is compact: \texttt{<correlationId>|<lastLine>}. A replica may return an empty last line if its file is empty, and the client treats the first arriving valid reply as sufficient.

\subsection{READ\_ALL Scenario}

The consistency-oriented path returns every local line, followed by an end marker.

\begin{lstlisting}
private void handleReadAll(String correlationId, String replyQueueName) {
    List<String> allLines = FileUtils.readAllLines(dataFilePath);

    for (String line : allLines) {
        String response = correlationId + "|" + replicaId + "|" + line;
        rabbitTemplate.convertAndSend(replyQueueName, response);
    }

    String endMarker = correlationId + "|END|" + replicaId;
    rabbitTemplate.convertAndSend(replyQueueName, endMarker);
}
\end{lstlisting}

This is the key mechanism that enables client-side reconciliation. Every line is sent separately, and the client waits until it sees an \texttt{END} marker from each replica or until the timeout expires.

\section{Client Writer}

The writer is a profile-scoped command-line runner. It either sends one line from the command line or enters an interactive loop.

\begin{lstlisting}
@Override
public void run(String... args) throws Exception {
    List<String> inputArgs = applicationArguments.getNonOptionArgs();

    if (!inputArgs.isEmpty()) {
        String line = String.join(" ", inputArgs);
        sendLine(line);
    } else {
        Scanner scanner = new Scanner(System.in);
        while (true) {
            String line = scanner.nextLine();
            if ("exit".equalsIgnoreCase(line)) {
                break;
            }
            if (!line.trim().isEmpty()) {
                sendLine(line);
            }
        }
    }

    System.exit(0);
}
\end{lstlisting}

The actual send operation publishes to the fanout exchange:

\begin{lstlisting}
private void sendLine(String line) {
    rabbitTemplate.convertAndSend(RabbitMQConfig.WRITE_EXCHANGE, "", line);
    System.out.println("Sent: " + line);
}
\end{lstlisting}

Because the exchange is fanout, the routing key is not used. The broker broadcasts the same payload to every bound replica queue.

\section{Client Reader: Availability-First Read}

\texttt{ClientReader} implements the simpler read mode. It asks every replica for the last line, but it accepts the first reply that comes back.

\begin{lstlisting}
String correlationId = UUID.randomUUID().toString();
String replyQueueName = "reply_" + correlationId;

Queue replyQueue = new Queue(replyQueueName, false, false, true);
rabbitAdmin.declareQueue(replyQueue);

CountDownLatch latch = new CountDownLatch(1);
AtomicReference<String> result = new AtomicReference<>();
\end{lstlisting}

The client creates a temporary exclusive queue for replies. That queue exists only for the duration of the request and is deleted afterwards.

The listener waits for the first reply and stores it:

\begin{lstlisting}
container.setMessageListener(message -> {
    String reply = new String(message.getBody());
    result.set(reply);
    latch.countDown();
});
\end{lstlisting}

Then it sends the request:

\begin{lstlisting}
String request = "READ_LAST:" + correlationId + ":" + replyQueueName;
rabbitTemplate.convertAndSend(RabbitMQConfig.READ_REQUEST_EXCHANGE, "", request);
\end{lstlisting}

Finally, it waits up to the configured timeout and prints the returned last line if one arrives.

This mode models an availability-oriented decision: the client prefers a quick answer over a globally reconciled one. If multiple replicas disagree, the client does not try to resolve the conflict here.

\section{Client Reader V2: Majority Voting Read}

The second reader mode is more elaborate. It asks every replica for all of its lines and then performs a quorum-based merge.

\subsection{Data Collection}

The client tracks two structures: one for the lines received from each replica and one for the replicas that have sent their \texttt{END} marker.

\begin{lstlisting}
Map<String, Set<String>> replicaLines = new HashMap<>();
Map<String, Boolean> replicaEnded = new HashMap<>();
Object lock = new Object();
\end{lstlisting}

Each message is parsed into three parts. A data message has the form \texttt{correlationId|replicaId|line}. An end marker has the form \texttt{correlationId|END|replicaId}.

\begin{lstlisting}
String[] parts = reply.split("\\|", 3);
String corrId = parts[0];
String replicaId = parts[1];
String content = parts[2];

if ("END".equals(replicaId)) {
    replicaEnded.put(content, true);
} else {
    replicaLines.computeIfAbsent(replicaId, k -> new HashSet<>()).add(content);
}
\end{lstlisting}

The code ignores replies whose correlation identifier does not match the current request. That matters because temporary reply queues and asynchronous delivery can otherwise cause stale messages to leak into a new read.

\subsection{Waiting for Replica Completion}

The client waits until every replica has sent an end marker or the timeout expires.

\begin{lstlisting}
while (System.currentTimeMillis() - startTime < readAllTimeoutMs) {
    synchronized (lock) {
        if (replicaEnded.size() >= replicaCount) {
            break;
        }
    }
    Thread.sleep(100);
}
\end{lstlisting}

This makes the behavior deterministic enough for the demo while still tolerating partial failures. If a replica is down, the timeout will expire and the client will reconcile whatever data it has received.

\subsection{Majority Voting}

The reconciliation logic counts how many replicas produced each exact line, then keeps only the values that appear in a majority of replicas.

\begin{lstlisting}
private void applyMajorityVoting(Map<String, Set<String>> replicaLines) {
    Map<String, Integer> lineVotes = new HashMap<>();
    for (String replicaId : replicaLines.keySet()) {
        for (String line : replicaLines.get(replicaId)) {
            lineVotes.put(line, lineVotes.getOrDefault(line, 0) + 1);
        }
    }

    Map<Integer, String> finalLines = new TreeMap<>();
    int requiredVotes = (replicaCount / 2) + 1;
    for (String line : lineVotes.keySet()) {
        int votes = lineVotes.get(line);
        if (votes >= requiredVotes) {
            String[] parts = line.split(" ", 2);
            int lineNum = Integer.parseInt(parts[0]);
            finalLines.put(lineNum, line);
        }
    }
}
\end{lstlisting}

The majority threshold is computed as \texttt{(replicaCount / 2) + 1}. With three replicas, the threshold is two votes. This means a line must be present in at least two replicas to survive reconciliation.

The implementation assumes the first token of each line is a numeric prefix, such as \texttt{1 First line}. This is how the final output is sorted by line number. If a line cannot be parsed, it is skipped with a warning. That is a pragmatic choice for the assignment because the demo format naturally uses numbered records.

\section{Message Formats}

The codebase uses explicit string formats instead of serializing structured objects. This keeps the messaging easy to inspect during debugging.

\begin{longtable}{|p{0.28\textwidth}|p{0.62\textwidth}|}
\hline
\textbf{Purpose} & \textbf{Format} \\
\hline
Write request & Any plain line of text, often \texttt{<number> <content>} \\
\hline
READ\_LAST request & \texttt{READ\_LAST:<correlationId>:<replyQueueName>} \\
\hline
READ\_LAST response & \texttt{<correlationId>|<lastLine>} \\
\hline
READ\_ALL request & \texttt{READ\_ALL:<correlationId>:<replyQueueName>} \\
\hline
READ\_ALL data response & \texttt{<correlationId>|<replicaId>|<line>} \\
\hline
READ\_ALL end marker & \texttt{<correlationId>|END|<replicaId>} \\
\hline
\end{longtable}

These conventions are the glue that lets the reader clients interpret asynchronous replies safely.

\section{Scenarios and What They Demonstrate}

\subsection{All Replicas Healthy}

When all replicas are running, every write is appended to all three local files. In this case, \texttt{reader} and \texttt{readerv2} usually return the same logical result because the replicas agree.

\subsection{Replica Down During Writes}

If one replica is offline, it misses the write messages published while it is down. When it comes back later, it does not automatically recover missing lines. This creates a visible divergence between the files.

This scenario is essential to the project because it shows why replication alone does not guarantee consistency. Delivery failures and downtime create disagreement unless there is a repair protocol.

\subsection{READ\_LAST as an Availability-First Choice}

The simple reader chooses the first replica that replies. This gives a fast answer and keeps the read path resilient to slow or failing replicas, but it does not guarantee that the answer reflects the majority state.

\subsection{READ\_ALL with Majority Voting}

The second reader gathers all visible lines and then applies quorum logic. In the normal three-replica setup, a line survives only if at least two replicas agree on it.

This is the code path that best illustrates the CAP trade-off in the project. It sacrifices simplicity and latency for a stronger reconciliation result.

\subsection{Empty Files and Missing Data}

The helper methods return empty values rather than throwing exceptions when files are missing or empty. That makes the entire system more robust during startup and after a fresh deployment.

\subsection{Malformed or Conflicting Lines}

The reconciliation code expects lines to begin with a numeric prefix. If that assumption is violated, the line is skipped. This is not a general-purpose merge algorithm; it is a project-specific reconciliation strategy tailored to the expected input format.

\section{What Each Class Is Responsible For}

\begin{itemize}[leftmargin=1.4em]
  \item \textbf{TpCapReplicationApplication}: starts Spring Boot.
  \item \textbf{RabbitMQConfig}: declares exchanges, queues, and bindings.
  \item \textbf{FileUtils}: handles directory creation and file reads and writes.
  \item \textbf{ReplicaService}: persists writes and answers read requests.
  \item \textbf{ClientWriter}: publishes writes to the broker.
  \item \textbf{ClientReader}: reads one last line from the first responsive replica.
  \item \textbf{ClientReaderV2}: aggregates all replica data and reconciles it by majority.
\end{itemize}

\section{Conclusion}

The project is a compact but instructive replication example. It combines Spring Boot profiles, RabbitMQ fanout exchanges, temporary reply queues, and file-based persistence to demonstrate how a distributed system can prefer either quick availability or a more consistent merged view. The code is intentionally direct, which makes it well suited for teaching, debugging, and documenting the basic mechanics of replication under partial failure.

\end{document}